\paragraph{问题陈述.}
激光超声（LUS）能够对关键基础设施进行非接触式检测；然而，采集的信号不可避免地受到强烈环境噪声和光子噪声的污染，信噪比（SNR）可低至$-10$~dB\cite{laser_ultrasound_noise}。这种低SNR从根本上限制了无损检测（NDT）中缺陷检测的可靠性，表现为三种具体形式：到达时间抖动降低了厚度和速度测量精度\cite{arrival_time_jitter}，振幅失真掩盖了疲劳裂纹和腐蚀\cite{amplitude_distortion}，以及升高的假阴性率导致危险缺陷未被检测到\cite{false_negative_ndt}。传统滤波方法——带通滤波、小波去噪和谱减法——以固定先验进行操作，无法适应激光产生声发射特有的非平稳、结构化噪声\cite{conventional_filter_limits}。虽然深度学习（DL）在声学成像任务的去噪中表现出强大的能力\cite{deep_learning_ultrasound}，但现有在加性高斯白噪声（AWGN）上训练的DL架构存在严重的仿真到实际差距：它们在部署于结构化实验室或现场噪声时过度平滑波形形态、
模糊到达边缘并抑制合法信号成分\cite{over_smoothing_issue,sim_to_real_gap}。在实现稳健SNR改善的同时保持波形保真度仍然是激光超声无损检测中一项开放的研究挑战。